\subsection{Submersions and quotient varieties}

5.9.1. Let \(f: X \to Y\) be a morphism of varieties, \(a\) a point of \(X\) and let us set \(b = f(a)\) . The following conditions are equivalent:

\begin{enumerate}
\def\labelenumi{(\roman{enumi})}
\item
  The linear map \(T_{a}(f)\) is surjective, and its kernel admits a topological complement in \(T_{a}(X)\) .
\item
  There exists a chart (U, \(\varphi\) , E) of X at \(a\) , a chart (V, \(\psi\) , F) of Y at \(b\) and a surjective continuous linear map \(u\) from E to F such that
\end{enumerate}

\[
f (\mathrm{U}) \subset \mathrm{V}, \quad \psi \circ f = u \circ \varphi
\]

and that the kernel of \(u\) admits a topological supplement in E.

\begin{enumerate}
\def\labelenumi{(\roman{enumi})}
\setcounter{enumi}{2}
\item
  There exist an open neighborhood U of a, an open neighborhood V of b containing \(f(\mathbf{U})\) and a morphism g from U into a variety Z such that the map \((f,g)\) from U into \(V \times Z\) is an isomorphism.
\item
  There exists an open neighborhood V of b and a morphism s from V to X such that \(s(b) = a\) and \(f(s(y)) = y\) for all y in V (“local section”).
\end{enumerate}

When X and Y are finite-dimensional, the preceding conditions are equivalent to the following condition:

\begin{enumerate}
\def\labelenumi{(\alph{enumi})}
\setcounter{enumi}{21}
\tightlist
\item
  There exist an open neighborhood U of a, an open neighborhood V of b containing \(f(\mathbf{U})\) , and a coordinate system \((\eta^{1},\ldots,\eta^{n})\) on V such that the functions \(\eta^{i}\circ f\) on U be part of a coordinate system on U.
\end{enumerate}

If properties (i) through (iv) are satisfied, one says that \(f\) is a submersion at \(a\) . The set of points where \(f\) is a submersion is open in \(X\) ; if this open set is all of X, then f is said to be a submersion.

5.9.2. Let \(f: X \to Y\) and \(g: Y \to Z\) be two morphisms. If \(f\) and \(g\) are submersions, the same holds for \(g \circ f\) . Conversely, if \(g \circ f\) is a submersion and if \(f\) is surjective, then \(g\) is a submersion.

5.9.3. If \(f: X \to Y\) and \(f': X' \to Y'\) are submersions, \(f \times f'\) is a submersion.

5.9.4. A submersion \(f: X \to Y\) is an open mapping (cf. Top. Gen., chap. I, 4 \(^{e}\) ed., § 5, no. 1); in particular, the equivalence relation R defined by \(f\) is open, \(f\) defines, by passing to the quotient, a homeomorphism from X/R onto \(f(X)\) , and \(f(X)\) is open in Y.

5.9.5. Let R be an equivalence relation on a variety X. We say that R is regular if there exists on the quotient space X/R a variety structure such that the canonical projection \(p: X \to X/R\) is a submersion; this variety structure is then unique; it is a quotient of that of X (Sets, ch.~IV, § 2, no. 6): in other words, for a morphism \(g\) from X/R into a variety Y to exist, it is necessary and sufficient that \(g \circ p\) be a morphism from X to Y.

Let \(C \subset X \times X\) be the graph of R. For R to be regular, it is necessary and sufficient that the following two conditions be satisfied:

\begin{enumerate}
\def\labelenumi{(\roman{enumi})}
\item
  C is a submanifold of X × X.
\item
  The map \(pr_{1}\) from C into X is a submersion.
\end{enumerate}

(ii) also means that if \(a\) and \(b\) are congruent modulo \(\mathbf{R}\) there exists an open neighborhood U of \(a\) and a morphism \(s\) from U into X such that \(s(a) = b\) and \(s(x)\) is congruent to \(x\) modulo \(\mathbf{R}\) for all \(x \in \mathbf{U}\) .

Suppose R is regular. For the quotient variety X/R to be separated, it is necessary and sufficient that the graph of R be closed in \(X \times X\) .

5.9.6. Let X and X' be two varieties, R and R' regular equivalence relations on X and X', and let \(f: X \to X'\) be a morphism compatible with the relations R and R'. The map \(\bar{f}: X/R \to X'/R'\) deduced from \(f\) by passing to quotients is then a morphism.

5.9.7. (“Transitivity of quotients”) Let R and S be two equivalence relations on a variety X such that R implies S, and let S/R be the quotient equivalence relation on X/R. Suppose that R is regular. Then, for S to be regular, it is necessary and sufficient that S/R be regular; if this is the case, the canonical bijection

\[
(\mathrm{X} / \mathrm{R}) / (\mathrm{S} / \mathrm{R}) \rightarrow \mathrm{X} / \mathrm{S}
\]

is an isomorphism of varieties.

5.9.8. ("Products of quotients") Let \((\mathbf{X}_{i})_{i\in I}\) be a finite family of varieties, each equipped with a regular equivalence relation \(R_{i}\) . Let \(X = \prod_{i\in I} X_{i}\) and let R be the equivalence relation on X that is the product of the \(R_{i}\) (cf. Top. Gen., chap. I, 4th ed., § 5, no. 3, cor. of prop. 8). Then R is regular, and the canonical bijection from X/R onto \(\prod_{i\in I}(X_{i}/R_{i})\) is an isomorphism of varieties.

